% This is part of Mes notes de mathématique
% Copyright (c) 2011-2020
%   Laurent Claessens
% See the file fdl-1.3.txt for copying conditions.

%+++++++++++++++++++++++++++++++++++++++++++++++++++++++++++++++++++++++++++++++++++++++++++++++++++++++++++++++++++++++++++
\section{Groupes abéliens finis}
%+++++++++++++++++++++++++++++++++++++++++++++++++++++++++++++++++++++++++++++++++++++++++++++++++++++++++++++++++++++++++++

Source : \cite{FabricegPSFinis}.

Nous rappelons que l'exposant\index{exposant} d'un groupe fini est le \( \ppcm\) des ordres de ses éléments. Dans le cas des groupes abéliens finis, l'exposant joue un rôle important du fait qu'il existe un élément dont l'ordre est l'exposant. Cela est le théorème suivant.

\begin{theorem}[Exposant dans un groupe abélien fini]
    Un groupe abélien fini contient un élément dont l'ordre est l'exposant du groupe.
\end{theorem}

\begin{proof}
    Soit \( G\) un groupe abélien fini et \( x\in G\), un élément d'ordre maximum \( m\). Nous montrons par l'absurde que l'ordre de tous les éléments de \( G\) divise \( m\). Soit donc \( y\in G\), un élément dont l'ordre ne divise pas \( m\); nous notons $q$ son ordre. Vu que \( q\) ne divise pas \( m\), le nombre \( q\) possède au moins un facteur premier plus de fois que \( m\) : soit \( p\) premier tel que la décomposition de \( q\) contienne \( p^{\beta}\) et celle de \( m\) contienne \( p^{\alpha}\) avec \( \beta>\alpha\). Autrement dit,
    \begin{subequations}
        \begin{align}
            m=p^{\alpha}m'\\
            q=p^{\beta}q'
        \end{align}
    \end{subequations}
    où \( m'\) et \( q'\) ne contiennent plus le facteur \( p\). L'élément \( x\) étant d'ordre \( m\), l'élément \( x^{p^{\alpha}}\) est d'ordre \( m'\). De la même manière, l'élément \( y^{q'}\) est d'ordre \( p^{\beta}\). Étant donné que \( p^{\beta}\) et \( m'\) sont premiers entre eux, l'élément  \( x^{p^{\alpha}}y^{q'}\) est d'ordre \( p^{\alpha}m'>m\). D'où une contradiction avec le fait que \( x\) était d'ordre maximal.

    Par conséquent l'ordre de tous les éléments de $G$ divise celui de \( x\) qui est alors le \( \ppcm\) des ordres de tous les éléments de \( G\), c'est-à-dire l'exposant de \( G\).
\end{proof}

\begin{proposition} \label{PropfPRVxi}
    Soit \( G\) un groupe abélien fini et \( x\in G\), un élément d'ordre maximum. Alors
    \begin{enumerate}
        \item
            Il existe un morphisme \( \varphi\colon G\to \gr(x)\) tel que \( \varphi(x)=x\).
        \item   \label{ItemKRYwjU}
            Il existe un sous-groupe \( K\) de \( G\) tel que \( G=\gr(x)\oplus K\).
    \end{enumerate}
\end{proposition}

\begin{proof}
    Nous notons \( a\) l'ordre de \( x\) qui est également l'exposant du groupe \( G\).

    Nous allons prouver la première partie par récurrence sur l'ordre du groupe. Si \( G=\gr(x)\), alors c'est évident. Soit \( H\) un sous-groupe propre de \( G\) contenant \( x\) et tel que le problème soit déjà résolu pour \( H\) : il existe un morphisme \( \varphi\colon H\to \gr(x)\) tel que \( \varphi(x)=x\). Soit \( y\in G\setminus H\), d'ordre \( b\). Nous allons trouver un morphisme $\hat\varphi\colon \gr(H,y)\to \gr(x) $ telle que \( \hat\varphi(x)=x\).

    Pour cela nous commençons par construire les applications suivantes :
    \begin{equation}
        \begin{aligned}
            \tilde \varphi\colon \eZ/b\eZ\times H&\to \gr(x) \\
            (\bar k,h)&\mapsto x^{kl}\varphi(h)
        \end{aligned}
    \end{equation}
    où \( l\) est encore à déterminer, et
    \begin{equation}
        \begin{aligned}
            p\colon \eZ/b\eZ\times H&\to \gr(y,H) \\
            (\bar k,h)&\mapsto y^kh.
        \end{aligned}
    \end{equation}
    Pour que \( \tilde \varphi\) soit bien définie, il faut que \( a\) divise \( bl\). L'application \( p\) est bien définie parce que \( \bar k\) est pris dans \( \eZ/b\eZ\) et que \( b\) est l'ordre de \( y\).

    Nous allons construire le morphisme \( \hat \varphi\) en considérant le diagramme
    \begin{equation}
    \xymatrix{%
    \ker(p) \ar@{^{(}->}[r]        &   \eZ/b\eZ\times H\ar[d]_{\tilde \varphi}\ar[r]^p&\gr(y,H)\ar[ld]^{\hat \varphi}\\
          &   \gr(x)
       }
    \end{equation}
    que l'on voudra être commutatif. Vu que \( p\) est surjective, les théorèmes d'isomorphismes nous disent que
    \begin{equation}
        \gr(y,H)\simeq\frac{ \eZ/b\eZ\times H }{ \ker p }.
    \end{equation}
    Si \( [\bar k,h]\) est la classe de \( (\bar k,h)\) modulo \( \ker(p)\) alors nous voudrions définir \( \hat \varphi\) par
    \begin{equation}        \label{EqeesVxc}
        \hat\varphi\big( [\bar k,h] \big)=\tilde \varphi(\bar k,h).
    \end{equation}
    Pour que cela soit bien définit, il faut que si \( (\bar r,z)\in \ker p\),
    \begin{equation}
        \hat\varphi\big( [\bar k\bar r,hz] \big)=\hat\varphi\big( [\bar k,h] \big),
    \end{equation}
    c'est-à-dire que \( \tilde \varphi(\bar r,z)=e\). Du coup la définition \eqref{EqeesVxc} n'est bonne que si et seulement si
    \begin{equation}
        \ker(p)\subset\ker(\tilde\varphi ).
    \end{equation}
    Nous pouvons obtenir cela en choisissant bien \( l\).

    Déterminons d'abord le noyau de \( p\). Pour cela nous considérons un nombre \( \beta\) divisant \( b\) tel que \( \gr(y)\cap H=\gr(y^{\beta})\). Nous aurons \( p(\bar k,h)=e\) si et seulement si \( y^h=e\). En particulier \( h=y^{-k}\in\gr(y)\cap H=\gr(y^{\beta})\). Si \( h=(y^{\beta})^m=y^{m\beta}\), alors \( k=-m\beta\) et nous avons
    \begin{equation}
        \ker(p)=\{ (-m\beta,y^{m\beta})\tq m\in \eZ \}.
    \end{equation}
    En plus court : \( \ker(p)=\gr(\beta,y^{-\beta})\). Nous devons donc fixer \( l\) de telle sorte que \( \tilde \varphi(\beta,y^{-\beta})=e\). Étant donné que \( \varphi\) prend ses valeurs dans \( \gr(x)\), il existe un entier \( \alpha\) tel que \( \varphi(y^{-\beta})=x^{\alpha}\); en utilisant cet \( \alpha\), nous écrivons
    \begin{equation}
        \tilde \varphi(\beta,y^{-\beta})=x^{\beta l}\varphi(y^{-\beta})=x^{\beta l+\alpha}.
    \end{equation}
    Par conséquent nous choisissons \( l=-\alpha/\beta\). Nous devons maintenant vérifier que ce choix est légitime, c'est-à-dire que \( a\) divise \( bl\) et que \( \alpha/\beta\) est un entier.

    Étant donné que \( y\) est d'ordre \( b\),
    \begin{equation}
        e=\varphi(y^b)=\varphi(y^{-\beta b/\beta})=\varphi(y^{-\beta})^{b/\beta}=x^{b\beta/\alpha}.
    \end{equation}
    Par conséquent \( a\) divise \( \frac{ b\alpha }{ \beta }=-bl\).

    Pour voir que \( l\) est entier, nous nous rappelons que \( a\) est l'exposant de \( G\) (parce que \( x\) est d'ordre maximum) et que par conséquent \( b\) divise \( a\). Mais \( a\) divise \( \alpha\frac{ b }{ \beta }\). Donc \( \alpha/\beta\) est entier.

    Nous passons maintenant à la seconde partie de la preuve. Nous considérons un morphisme \( \varphi\colon G\to \gr(x)\) tel que \( \varphi(x)=x\). La première partie nous en assure l'existence. Nous montrons que
    \begin{equation}
        \begin{aligned}
            \psi\colon G&\to \gr(x)\oplus \ker(\varphi) \\
            g&\mapsto \big( \varphi(g),g\varphi(g)^{-1} \big)
        \end{aligned}
    \end{equation}
    est un isomorphisme. D'abord \( g\varphi(g)^{-1}\) est dans le noyau de \( \varphi\) parce que \( \varphi(g)^{-1}\) étant dans \( \gr(x)\), et \( \varphi\) étant un morphisme,
    \begin{equation}
        \varphi\big( g\varphi(g)^{-1} \big)=\varphi(g)\varphi(g)^{-1}=e.
    \end{equation}
    L'application \( \psi\) est un morphisme parce que, en utilisant le fait que \( G\) est abélien,
    \begin{subequations}
        \begin{align}
            \psi(g_1g_2)&=\big( \varphi(g_1g_2),g_1g_2\varphi(g_1g_2)^{-1} \big)\\
            &=\big( \varphi(g_1)\varphi(g_2),g_1\varphi(g_1)^{-1}g_2\varphi(g_2)^{-1} \big)\\
            &=\psi(g_1)\psi(g_2).
        \end{align}
    \end{subequations}
    L'application \( \psi\) est injective parce que si \( \psi(g)=(e,e)\) alors \( \varphi(g)=e\) et \( g\varphi(g)^{-1}=e\), ce qui implique \( g=e\).

    Enfin \( \psi\) est surjective parce qu'elle est injective et que les ensembles de départ et d'arrivée ont même cardinal. En effet par le premier théorème d'isomorphisme (théorème~\ref{ThoPremierthoisomo}) appliqué à \( \varphi\) nous avons
    \begin{equation}
        | G |=| \gr(x) |\cdot | \ker(\varphi) |.
    \end{equation}
\end{proof}

\begin{theorem} \label{ThoRJWVJd}
    Tout groupe abélien fini (non trivial) se décompose en
    \begin{equation}
        G\simeq \eZ/d_1\eZ\oplus\ldots\oplus \eZ/d_r\eZ
    \end{equation}
    avec \( d_1\geq 1\) et \( d_i\) divise \( d_{i+1}\) pour tout \( i=1,\ldots, r-1\).

    De plus la liste \( (d_1,\ldots, d_r)\) vérifiant ces propriétés est unique.
\end{theorem}

\begin{proof}
    Soit \( x_1\) un élément d'ordre maximal dans \( G\). Soit \( n_1\) son ordre et
    \begin{equation}
        H_1=\gr(x_1)=\eF_{n_1}.
    \end{equation}
    D'après la proposition~\ref{PropfPRVxi}\ref{ItemKRYwjU}, il existe un supplémentaire \( K_1\) tel que \( G=\eF_{n_1}\oplus K_1\). Si \( K_1=\{ 1 \}\) on s'arrête et on garde \( G=\eF_{n_1}\). Sinon on continue de la sorte en prenant \( x_2\) d'ordre maximal dans \( K_1\) etc.

    Nous devons maintenant prouver l'unicité de cette décomposition. Soit
    \begin{equation}
        G=\eF_{d_1}\oplus\ldots\oplus \eF_{d_r}=\eF_{s_1}\oplus\ldots\oplus \eF_{s_q}.
    \end{equation}
    L'exposant de \( G\) est \( d_r\) et \( s_q\). Donc \( d_r=s_q\). Les complémentaires étant égaux nous avons
    \begin{equation}
        \eF_{d_1}\oplus\ldots\oplus \eF_{d_{r-1}}=\eF_{s_1}\oplus\ldots\oplus \eF_{s_{q-1}}.
    \end{equation}
    En continuant nous trouvons \( r=q\) et \( d_i=s_i\).
\end{proof}

%+++++++++++++++++++++++++++++++++++++++++++++++++++++++++++++++++++++++++++++++++++++++++++++++++++++++++++++++++++++++++++
\section{Groupe monogène}
%+++++++++++++++++++++++++++++++++++++++++++++++++++++++++++++++++++++++++++++++++++++++++++++++++++++++++++++++++++++++++++
\label{SECooXIHPooWVSjhT}

Le théorème suivant donne quelques informations à propos de groupes monogènes. Il impliquera dans le corollaire~\ref{CORooMBLSooMHKmAq} qu'un groupe monogène d'ordre \( n\) possède \( \varphi(n)\) générateur où \( \varphi\) est la fonction indicatrice d'Euler définie en~\ref{DEFooWYIGooRVBTil}.

\begin{theorem}     \label{THOooDOMZooOEYHAe}
    Un groupe monogène est abélien. Plus précisément,
    \begin{enumerate}
        \item
            un groupe monogène infini est isomorphe à \( \eZ\),
        \item
            un groupe monogène fini est isomorphe à \( \eZ/n\eZ\) pour un certain \( n\).
    \end{enumerate}
\end{theorem}

\begin{proof}
    Le groupe est abélien parce que $g=a^n$, \( g'=a^{n'}\) implique \( gg'=q^{n+n'}=g'g\). Nous considérons un générateur \( a\) de \( G\) (qui existe parce que $G$ est monogène) et le morphisme surjectif
    \begin{equation}
        \begin{aligned}
            f\colon \eZ&\to G \\
            p&\mapsto a^p.
        \end{aligned}
    \end{equation}
    Si \( G\) est infini, alors \( f\) est injective parce que si \( a^n=a^{n'}\), alors \( a^{n-n'}=e\), ce qui rendrait \( G\) cyclique et par conséquent non infini. Nous concluons que si \( G\) est infini, alors \( f\) est une bijection et donc un isomorphisme \( \eZ\simeq G\).

    Si \( G\) est fini, alors \( f\) n'est pas injective et a un noyau \( \ker f\). Étant donné que \( \ker f\) est un sous-groupe de \( G\), il existe un (unique) \( n\) tel que \( \ker f=n\eZ\) et le premier théorème d'isomorphisme (théorème~\ref{ThoPremierthoisomo}) nous indique que
    \begin{equation}
        \eZ/\ker f=\eZ/n\eZ=\Image f=G.
    \end{equation}

\end{proof}

Le lemme suivant donne une démonstration alternative, avec une construction plus explicite de l'isomorphisme.

\begin{lemma}[\cite{MonCerveau}]   \label{LemZhxMit}

    À propos de groupes monogènes\footnote{Définition~\ref{DEFooWMFVooLDqVxR}.}

    \begin{enumerate}
        \item


    Soit un groupe monogène \( G\) d'ordre fini \( n\) dont \( g\) est un générateur. Alors il existe un isomorphisme
    \begin{equation}
        \phi\colon G\to (\eZ/n\eZ,+)
    \end{equation}
    tel que \( \phi(g)=1\).

\item

    Si \( G\) est un groupe monogène d'ordre infini et si \( g\) est un générateur, alors il existe un isomorphisme
    \begin{equation}
        \phi\colon G\to (\eZ,+)
    \end{equation}
    tel que \( \phi(g)=1\).

   \item

    Soient \( G\) et \( H\) deux groupes monogènes de même ordre. Soient \( g\) un générateur de \( G\) et \( h\), un générateur de \( H\). Il existe un isomorphisme de \( G\) sur \( H\) qui envoie \( g\) sur \( h\).
    \end{enumerate}
\end{lemma}

\begin{proof}

    Commençons par enfoncer une porte ouverte : vu que le groupe est monogène, l'ordre du groupe est égal à l'ordre de son générateur. Nous séparons les cas selon quel l'ordre soit fini ou non.

    \begin{subproof}
        \item[L'ordre de \( G\) est fini et vaut \( n\)]
            Si \( k\in\eZ\), nous notons \( [k]_n\) la classe de \( k\) modulo \( n\), c'est-à-dire l'ensemble \( \{ k+pn\tq p\in \eZ \}\).

            Nous construisons l'isomorphisme \( \phi\colon G\to \eZ/n\eZ\) de la façon suivante:
            \begin{equation}
                \phi(g^m)=[m]_n.
            \end{equation}
            Cela est une bonne définition parce que une égalité du type \( g^m=g^{m'}\) implique que \( m\) et \( m'\) soient dans la même classe modulo \( n\). Nous vérifions que cela est une isomorphisme entre \( G\) et \( \eZ/n\eZ\).

            \begin{subproof}
            \item[Morphisme]
                Pour l'identité, si \( x=e\) alors \( m=0\) et \( \phi(e)=[0]_n\). Et si \( x=g^k\), \( y=g^l\) alors \( \phi(xy)=\phi(g^{k+l})=[k+l]_n=[k]_n+[l]_n=\phi(x)+\phi(y) \).
            \item[Injectif]
                Supposons \( \phi(g^k)=\phi(g^l)\) avec \( k\geq l\). Nous avons \( h^k=h^l\), dont \( h^{k-l}=e\), ce qui donne \( k-l\in [0]_n\) ou encore \( [k]_n=[l]_n\). En particulier \( g^k=g^l\).
            \item[Surjectif]
                La classe \( [k]_n\) est l'image de \( g^k\).
            \end{subproof}

        \item[L'ordre de \( G\) est infini]

                Si l'ordre de \( G\) est infini alors un élément \( x\in G\) s'écrit de façon unique sous la forme \( x=g^m\) avec \( m\in \eZ\). Dans ce cas nous définissons directement \( \phi(g^m)=m\).

                Le reste de la preuve est alors identique au cas d'ordre fini, mais sans les complications liées au modulo.

    \end{subproof}

    La dernière assertion s'obtient des précédentes par composition d'isomorphismes.

\end{proof}


\begin{corollary}\label{CORooMBLSooMHKmAq}
    Un groupe monogène d'ordre \( n\) possède \( \varphi(n)\) générateurs où \( \varphi\) est la fonction indicatrice d'Euler définie en~\ref{DEFooWYIGooRVBTil}.
\end{corollary}

\begin{proof}
    Le théorème~\ref{THOooDOMZooOEYHAe} nous dit qu'un groupe monogène d'ordre \( n\) est isomorphe à \( \eZ/n\eZ\). La proposition~\ref{PropZnmuphiGensn} nous indique que \( \eZ/n\eZ\) possède \( \varphi(n)\) générateurs.
\end{proof}

%+++++++++++++++++++++++++++++++++++++++++++++++++++++++++++++++++++++++++++++++++++++++++++++++++++++++++++++++++++++++++++
\section{Automorphismes du groupe \texorpdfstring{$ \eZ/n\eZ$}{Z/nZ}}
%+++++++++++++++++++++++++++++++++++++++++++++++++++++++++++++++++++++++++++++++++++++++++++++++++++++++++++++++++++++++++++

Notons que \( \eZ/n\eZ=\eF_n\) est un groupe pour l'addition tandis que \( (\eZ/n\eZ)^*\) est un groupe pour la multiplication. Il ne peut donc pas y avoir d'équivoque.

\begin{theorem}[\cite{ooUDKTooTWYpzN}]   \label{ThoozyeSn}
    Pour chaque \( x\in (\eZ/n\eZ)^*\) nous considérons l'application
    \begin{equation}
        \begin{aligned}
            \sigma_x\colon \eZ/n\eZ&\to \eZ/n\eZ \\
            y&\mapsto xy.
        \end{aligned}
    \end{equation}
    L'application
    \begin{equation}
        \sigma\colon \big( (\eZ/n\eZ)^*,\cdot\big)\to \Aut\big( \eZ/n\eZ,+ \big)
    \end{equation}
    ainsi définie est un isomorphisme de groupes.
\end{theorem}
L'énoncé de ce théorème s'écrit souvent rapidement par
\begin{equation}
    \Aut(\eZ/n\eZ)=(\eZ/n\eZ)^*,
\end{equation}
mais il faut bien garder à l'esprit qu'à gauche on considère le groupe additif et à droite celui multiplicatif.

\begin{proof}
    Nous notons \( [x]\) la classe de \( x\) dans \( \eZ/n\eZ\). Nous avons \( \eZ/n\eZ=[1]\). Soit \( f\) un automorphisme de \( (\eZ/n\eZ,+)\); pour tout \( r\in \eZ\) nous avons
    \begin{equation}
        f([r])=f(r[1])=rf([1])=[r]f([1]).
    \end{equation}
En particulier, vu que \( f\) est surjective, il existe un \( r\) tel que \( f([r])=[1]\). Pour un tel \( r\) nous avons \( [1]=[r]f([1])\), c'est-à-dire que nous avons montré que \( f([1])\) est inversible dans \(  \big( (\eZ/n\eZ)^*,\cdot\big)\). Nous montrons à présent que\footnote{Le \( \sigma\) donné ici est l'inverse de celui donné dans l'énoncé. Cela ne change évidemment rien à la validité de l'énoncé et de la preuve.}
    \begin{equation}
        \begin{aligned}
            \sigma\colon \Aut( (\eZ/n\eZ,+))&\to \big( (\eZ/n\eZ)^*,\cdot \big)\\
            f&\mapsto f([1])
        \end{aligned}
    \end{equation}
    est un isomorphisme.

    Nous commençons par la surjectivité. Soit \( [a]\in (\eZ/n\eZ)^*\). Les élément \( [a]\) et \( [1]\) étant tous deux des générateurs de \( (\eZ/n\eZ,+)\), il existe un automorphisme de \( \eZ/n\eZ\) qui envoie \( [1]\) sur \( [a]\) par le lemme~\ref{LemZhxMit}. Cela prouve la surjectivité de \( \sigma\).

    En ce qui concerne l'injectivité, considérons des automorphismes \( f_1\) et \( f_2\) de \( (\eZ/n\eZ,+)\) tels que \( f_1([1])=f_2([1])\). Les automorphismes \( f_1\) et \( f_2\) prennent la même valeur sur un générateur et donc sur tout le groupe. Donc \( f_1=f_2\).

    Enfin nous prouvons que \( \sigma\) est un morphisme, c'est-à-dire que \( \sigma(f\circ g)=\sigma(f)\sigma(g)\). Nous avons
    \begin{subequations}
        \begin{align}
            f\big( g([1]) \big)&=f\big( g([1])[1] \big)=g([1])f([1])=\sigma(f)\sigma(g).
        \end{align}
    \end{subequations}
\end{proof}

Ce dernier résultat s'étend aux groupes cycliques.
\begin{proposition}     \label{PROPooBZOMooVOHoYf}
    Si \( G\) est un groupe cyclique\footnote{Définition \ref{DefHFJWooFxkzCF}.} d'ordre \( n\), alors
    \begin{equation}
        \Aut(G)=(\eZ/n\eZ)^*.
    \end{equation}
\end{proposition}
